\documentclass[11pt,a4paper]{article}

\usepackage[margin=2.3cm]{geometry}
\usepackage{amsmath}
\usepackage{booktabs}
\usepackage{array}
\usepackage{enumitem}
\usepackage{hyperref}
\usepackage[T1]{fontenc}
\usepackage{lmodern}

\title{Milestone 2: Scientific Contextualization and Tradeoff Analysis\\
\large List-Based Scheduling for High-Level Synthesis}
\author{Student: Md Azadul Islam\\Student ID: 2220070}
\date{28 June 2026}

\begin{document}
\maketitle

\section*{Purpose of this Milestone}

This milestone places list-based scheduling into the wider scientific and engineering context of High-Level Synthesis (HLS).
The focus is not only to list related papers, but to compare approaches, assumptions, design alternatives, tradeoffs, and consequences for embedded systems and constrained hardware.

\section{Position of List-Based Scheduling in the HLS Flow}

\begin{itemize}[leftmargin=*]
    \item High-Level Synthesis converts a behavioral description, for example C/C++ or an algorithmic model, into Register-Transfer Level (RTL) hardware.
    \item A simplified HLS flow is:
    \[
    \text{Algorithm} \rightarrow \text{DFG/CDFG} \rightarrow \text{Scheduling} \rightarrow \text{Allocation} \rightarrow \text{Binding} \rightarrow \text{RTL}
    \]
    \item Scheduling decides \textbf{when} operations execute.
    \item Allocation decides \textbf{how many} hardware resources are available.
    \item Binding decides \textbf{which operation uses which resource}.
    \item List-based scheduling mainly addresses the scheduling step, but its result strongly affects allocation, binding, area, latency, and control complexity.
\end{itemize}

\section{Related Approaches and Neighboring Methods}

\begin{table}[h]
\centering
\small
\begin{tabular}{p{0.22\linewidth}p{0.31\linewidth}p{0.34\linewidth}}
\toprule
Approach & Main idea & Relation to list-based scheduling \\
\midrule
ASAP scheduling & Schedule every operation as early as possible. & Often used as a baseline and for calculating earliest times. \\
ALAP scheduling & Schedule every operation as late as possible within a latency bound. & Often used with ASAP to compute mobility. \\
List-based scheduling & Maintain a ready list and schedule operations according to priority while respecting resources. & Main topic; practical heuristic for resource-constrained scheduling. \\
Force-directed scheduling & Uses probability/distribution forces to balance operation placement over time. & More global than simple list scheduling, but usually more complex. \\
ILP/exact scheduling & Formulates scheduling as an optimization problem. & Can find optimal schedules for small cases, but may scale poorly. \\
Modulo/pipelined scheduling & Overlaps loop iterations to improve throughput. & Important for loops and streaming applications; different objective from simple acyclic DFG scheduling. \\
Manual RTL design & Designer explicitly controls cycles, resources, and registers. & Gives high control but needs more hardware expertise and time. \\
Commercial/open-source HLS tools & Automate scheduling, allocation, binding, and RTL generation. & List scheduling or variants can appear as internal scheduling heuristics. \\
\bottomrule
\end{tabular}
\end{table}

\section{Scientific Context}

\subsection*{Classical HLS Research}

\begin{itemize}[leftmargin=*]
    \item Classical HLS research studies how to synthesize efficient hardware from high-level descriptions.
    \item Scheduling, allocation, and binding are central problems in this area.
    \item De Micheli explains HLS and digital circuit synthesis as part of architectural and logic-level design optimization \cite{DeMicheli1994}.
    \item Gajski et al. discuss HLS foundations, design representations, scheduling, allocation, and system-level design issues \cite{Gajski1992}.
    \item Paulin and Knight introduced influential force-directed scheduling work for behavioral synthesis, which is important for comparison because it is more global than simple list scheduling \cite{PaulinKnight1989}.
\end{itemize}

\subsection*{Neighboring Problem Classes}

\begin{itemize}[leftmargin=*]
    \item \textbf{Compiler instruction scheduling:} also schedules operations under dependencies and resource constraints.
    \item \textbf{Real-time scheduling:} focuses on deadlines and task timing; HLS scheduling focuses more on operation-level hardware cycles.
    \item \textbf{Resource-constrained project scheduling:} similar graph/resource structure, but usually not hardware-specific.
    \item \textbf{Hardware/software co-design:} scheduling decisions may affect whether work is done in software, hardware, or accelerators.
    \item \textbf{Loop pipelining and dataflow HLS:} focuses on throughput, initiation interval, and streaming behavior.
\end{itemize}

\section{Main Design Alternatives}

\subsection*{Alternative 1: ASAP/ALAP Only}

\begin{itemize}[leftmargin=*]
    \item Very simple and easy to understand.
    \item Useful for estimating earliest/latest control steps.
    \item Not enough when resource constraints are strict.
    \item Example: ASAP may place two additions in the same cycle even if only one ALU exists.
\end{itemize}

\subsection*{Alternative 2: List-Based Scheduling}

\begin{itemize}[leftmargin=*]
    \item Practical compromise between simplicity and useful scheduling quality.
    \item Works step by step using a ready list.
    \item Can include different priority rules:
    \begin{itemize}
        \item critical path,
        \item mobility,
        \item number of successors,
        \item operation latency,
        \item resource urgency.
    \end{itemize}
    \item Main limitation: the result depends strongly on the priority rule.
\end{itemize}

\subsection*{Alternative 3: Force-Directed Scheduling}

\begin{itemize}[leftmargin=*]
    \item Tries to distribute operations over time to reduce resource peaks.
    \item More global reasoning than normal list scheduling.
    \item Useful when resource sharing and balanced hardware usage are important.
    \item More complex to explain and implement than simple list scheduling.
\end{itemize}

\subsection*{Alternative 4: Exact Optimization, for Example ILP}

\begin{itemize}[leftmargin=*]
    \item Scheduling can be formulated as an optimization problem.
    \item Can optimize objectives such as minimum latency or minimum resource usage.
    \item Strong advantage: can give optimal results for small examples.
    \item Main weakness: complexity increases quickly with graph size, resource types, and timing constraints.
\end{itemize}

\subsection*{Alternative 5: Loop Pipelining / Modulo Scheduling}

\begin{itemize}[leftmargin=*]
    \item Important for loops and streaming computations.
    \item Objective is often throughput, not only total latency.
    \item Key metric:
    \[
    II = \text{Initiation Interval}
    \]
    \item A low initiation interval means new loop iterations can start frequently.
    \item This is related to scheduling but should not be confused with simple list scheduling for an acyclic expression graph.
\end{itemize}

\section{Tradeoff Analysis}

\begin{table}[h]
\centering
\small
\begin{tabular}{p{0.25\linewidth}p{0.33\linewidth}p{0.33\linewidth}}
\toprule
Tradeoff & List-based scheduling advantage & Limitation / risk \\
\midrule
Optimality vs feasibility & Fast and feasible for practical graphs. & Does not guarantee globally optimal schedule. \\
Runtime vs schedule quality & Much faster than many exact methods. & May produce worse latency or resource usage than optimal methods. \\
Simplicity vs global reasoning & Easy to understand and implement. & Local priority decisions can cause later conflicts. \\
Area vs latency & More resources can reduce cycles. & More functional units increase area and power. \\
Determinism vs flexibility & Fixed priority rules give predictable behavior. & Changing priority rules can change the final schedule. \\
Offline computation vs runtime adaptation & HLS scheduling is normally done offline before hardware generation. & Less adaptable to runtime changes than software scheduling. \\
Control simplicity vs performance & Simple schedules often produce simpler controllers. & Aggressive schedules may require more complex control and registers. \\
Memory simplicity vs realism & Basic DFG examples ignore many memory conflicts. & Real HLS designs often depend strongly on memory ports and arrays. \\
\bottomrule
\end{tabular}
\end{table}

\section{Assumptions Behind Different Approaches}

\begin{itemize}[leftmargin=*]
    \item \textbf{ASAP/ALAP:}
    \begin{itemize}
        \item assumes dependency information is enough for initial timing estimation;
        \item may ignore detailed resource conflicts in the simplest form.
    \end{itemize}
    \item \textbf{List-based scheduling:}
    \begin{itemize}
        \item assumes operations can be ordered using a useful priority function;
        \item assumes resource availability is known;
        \item assumes local greedy decisions are acceptable.
    \end{itemize}
    \item \textbf{Force-directed scheduling:}
    \begin{itemize}
        \item assumes operation distributions and probabilities can guide balanced scheduling;
        \item works better when global resource balance is more important.
    \end{itemize}
    \item \textbf{ILP/exact methods:}
    \begin{itemize}
        \item assumes the scheduling problem can be encoded accurately;
        \item assumes available solver time is acceptable.
    \end{itemize}
    \item \textbf{Loop pipelining:}
    \begin{itemize}
        \item assumes repeated loop iterations and throughput-oriented optimization;
        \item requires attention to loop-carried dependencies and memory ports.
    \end{itemize}
\end{itemize}

\section{Consequences for Embedded Systems and Constrained Hardware}

\begin{itemize}[leftmargin=*]
    \item Embedded systems often have limited:
    \begin{itemize}
        \item chip area,
        \item power budget,
        \item memory bandwidth,
        \item timing margins,
        \item and design time.
    \end{itemize}
    \item List-based scheduling is attractive because it is practical and understandable.
    \item It can help show how adding one extra ALU changes latency.
    \item For small FPGA or microcontroller-support hardware, resource constraints are very important.
    \item For safety-critical or hard real-time embedded systems, deterministic timing and explainable design decisions are important.
    \item For high-throughput signal processing or image processing, loop pipelining and dataflow scheduling may be more important than simple acyclic list scheduling.
\end{itemize}

\section{Comparison with My Implementation Direction}

The topic can be connected to the implementation work using a small arithmetic expression:

\[
z = (a+b)\times(c+d)+e
\]

\begin{itemize}[leftmargin=*]
    \item With \textbf{1 ALU + 1 multiplier}, the two additions must be separated:
    \[
    4 \text{ cycles}
    \]
    \item With \textbf{2 ALUs + 1 multiplier}, the two additions can run in parallel:
    \[
    3 \text{ cycles}
    \]
    \item This demonstrates the tradeoff:
    \[
    \text{More hardware resources} \Rightarrow \text{lower latency but higher area/cost}
    \]
    \item This implementation example is useful because it makes the scheduling tradeoff visible in a schedule table, C++ output, or Questa/VHDL waveform.
\end{itemize}

\section{Why Some Related Work Is Only Indirectly Related}

\begin{itemize}[leftmargin=*]
    \item Some HLS papers focus mainly on modern tool flows, pragmas, or compiler infrastructure instead of the scheduling algorithm itself.
    \item Some FPGA papers focus on placement/routing or bitstream generation, which happens after RTL generation and is therefore less directly related.
    \item Some real-time scheduling papers focus on task deadlines rather than operation-level hardware scheduling.
    \item Some machine-learning accelerator papers use HLS, but their main contribution may be architecture or dataflow design, not list scheduling.
    \item Therefore, the final paper should use only references that help explain scheduling, resource constraints, or HLS tradeoffs directly.
\end{itemize}

\section{Selected References and Justification}

\begin{table}[h]
\centering
\small
\begin{tabular}{p{0.25\linewidth}p{0.62\linewidth}}
\toprule
Reference & Why it is relevant \\
\midrule
De Micheli \cite{DeMicheli1994} & Strong foundational source for digital synthesis, optimization, scheduling, and architectural-level design. \\
Gajski et al. \cite{Gajski1992} & Foundational HLS textbook; useful for explaining HLS flow and basic scheduling concepts. \\
Paulin and Knight \cite{PaulinKnight1989} & Important related method because force-directed scheduling is a major alternative to simple list scheduling. \\
Coussy and Morawiec \cite{Coussy2008} & Useful broader HLS context and bridge from algorithm to digital circuit. \\
Canis et al. \cite{Canis2011} & Useful modern tool-flow reference showing HLS relevance for FPGA processor/accelerator systems. \\
\bottomrule
\end{tabular}
\end{table}

\section{Excluded or Less Central References}

\begin{itemize}[leftmargin=*]
    \item Pure FPGA place-and-route references are less central because list scheduling happens earlier in the HLS flow.
    \item General operating-system scheduling references are less central because they schedule software tasks, not hardware operations in a DFG.
    \item Application-specific accelerator papers are useful only if they clearly discuss HLS scheduling or resource constraints.
    \item Vendor manuals can be useful for practical tools, but the seminar paper should mainly use scientific references for the core theory.
\end{itemize}

\section{Milestone 2 Checklist Status}

\begin{table}[h]
\centering
\begin{tabular}{p{0.74\linewidth}c}
\toprule
Requirement & Status \\
\midrule
Compared approaches or assumptions, not only summarized papers & Done \\
Identified neighboring approaches and design alternatives & Done \\
Explained why related work can be direct, indirect, or sparse & Done \\
Analyzed scientific and engineering tradeoffs & Done \\
Considered embedded-system constraints & Done \\
Justified selected references & Done \\
Justified excluded or less relevant references & Done \\
\bottomrule
\end{tabular}
\end{table}

\section*{Short Summary}

List-based scheduling is positioned between very simple ASAP/ALAP timing methods and more complex methods such as force-directed scheduling or exact optimization.
Its main advantage is practical feasibility and understandable scheduling behavior.
Its main weakness is that it is heuristic and priority-dependent.
For embedded systems, the most important tradeoff is usually latency versus area and power.

\bibliographystyle{plain}
\bibliography{references_milestone2_list_based_scheduling}

\end{document}
